\section{我的难点}

这一章记我反复要求、而子 agent 反复没做到的地方。执行链那张图讲的是这套东西怎么工作，不属于难点，放到第 2 章开篇去了。

\subsection{视觉一致：我每次只点名一处，它就只改一处}

界面上的不统一，每次都得我自己开页面对一遍，逐条说、逐条验。

招式卡高矮不一，原因是一张卡里描述长短不同，整排跟着参差；对局页左右两侧的格式也要一致，两侧标题长度不同，会把各自的筛选条推到不同的起点，看上去就不对称。我第一次说这件事时把方向指错了，它照着改了一遍间距，问题还在；真正管用的是给标题固定宽度。筛选按钮我要求对称、撑满本栏，与下面的伙伴网格同宽。还有「我方」两个字，写在分屏标题里既没用又不对齐，删掉之后两侧结构才真的一样；\code{X号位} 标记也是这样，改成绝对定位之后，已经选进队伍的卡不再因为多出一行而变高。

这些改动的代码量都很小，难的是我每次只能指出我看到的那一处，它就把那一处改掉，下一次我再看见另一处。

\subsection{重复信息：同一句话出现在几个地方}

同一句话出现在几个地方，只有真的用起来才会撞上。出征页顶部的徽标写着「训练 · PVE」或「对局 · PVP」，右栏又写了一遍同样的模式。回合结算后，横幅把刚在动画里念过、也已经写进战斗记录的那一帧复述一遍。最严重的一次是一段老师的讲解同时出现在顶部条、左下浮层和右下气泡三处，文字一模一样。

这类问题的梳理方式和上面完全不同：它不需要开页面逐处看，需要把玩家可见的全部文案集中到一张表上逐条对，「同一件事」才浮出来。我还漏过一类更小的——同一个动作在界面上有三个名字：标签写「认输」，按钮写「确认撤退」，战斗记录里又写成「撤退」。

\subsection{说了很多遍「做完了」，其实是没生效}

这一类的每次交付都配着一句「已完成」，而每一次都是单元测试全绿、功能没跑起来：我说过很多遍，仍然要一遍遍自己验。

军师改造时删掉一个变量的定义、留下它的引用，表现是控制台每秒抛一次错。敌方补位「决定了却无人提交」：换宠选好、存进变量就返回，玩家那条路径又拒绝替它提交，界面永久冻结在「正在出招」，那一局走不下去。陪练的主动气泡写完，全仓库没有调用点。最后一次是角色路由：我指名要陪练，回答却从老师那条路出来——\code{coach/client.js} 在发送前会追加一段请求说明，里面含有「回合」，而路由判断先匹配这个词，于是角色被改判。单元测试测的是函数，测不出「是不是它被调用了」。

现在每个能力都有一条能执行的入口检查（\code{wiring.test.js} 逐个断言），改完路由我必须自己在浏览器里跑一遍。光重申没有用，只能把它变成机械检查。

\subsection{检查永远通过}

比漏改更麻烦的是：用来兜底的检查本身就不起作用。

冒烟测试里的伙伴卡数量写死过，宠物增加之后那个数字不再对得上；语义检索读的语料是手工维护的第三份副本，没有生成器，两张卡改了它不知道，模型引用到的仍是旧文案；禁止词表检查只覆盖运行时构造得出来的分支，漏掉了一整条理由的正文；工具参数用合法行动列表的下标，而索引会随能量、道具和倒下整体位移，同一个下标在不同回合指向不同行动，回执里也没法核对。

这四条形态不同、毛病一样：检查写死了预期值，或者与被测对象不同源，或者只覆盖跑得到的那条路。所以每写完一条检查，我会先做一次负向验证——把错的东西放进去，看它会不会失败；再让期望值从被测对象自己推导出来，不从手写的常量来。
